\section{結言}
今回の特徴は、エージェントの価値が回答の品質から実際の業務やシステムへの介入能力へ移りつつある点にある。Shirudoでは稼働中のLinuxを防護し、Claude ShipではPRの修正と承認を担う。いずれも、人間が設定する制約や検証手続きの設計が実用化の前提となる。

研究領域でも、量子回路の形式検証や推薦システムの実験自動化が、探索と評価をつなぐ方向性を示した。対照的に、個人のローカル推論報告は未検証であり、イベント告知とともに観測窓外の投稿である。今後は各製品の実運用条件、再現可能な評価、エージェントの権限管理が継続的な注目点となる。
